% The DevOps track is grounded in CI/CD and GPU operations, without inventing IaC expertise.
\renewcommand{\cvfootertitle}{Cover Letter}
\renewcommand{\cvBody}{%
    \cvletterhead
        {Stuttgart, 25 October 2026}
        {Scaleway\\HDV Tribe -- Hardware Operations\\Paris, France}
        {Application: Backend Software Engineer (Python / DevOps)}

    Dear Scaleway team,

    I want Europe to build strong cloud and AI capabilities of its own and reduce its dependence on foreign providers. Your Hardware Operations role offers a concrete way to contribute to that ambition: building the software that keeps physical infrastructure reliable. QCS's position between application development and hardware performance closely matches the work I enjoy at Fraunhofer IPA.

    I have maintained several institute repositories in Python, Rust and C++, alongside responsibility for GitLab CI/CD, DGX A100/H200 GPU infrastructure and Grafana/Prometheus monitoring. BuildKit and sccache caching helped reduce builds from hours to minutes. This work has made maintenance, reproducibility and diagnosing failures part of my everyday engineering practice. I would bring that discipline to QCS's hardware-update, testing and inventory systems, with Python development as my core contribution and DevOps as my secondary strength.

    I am comfortable working where software meets hardware constraints. I retrofitted a Berghof PLC with a Hailo accelerator and deployed object detection on it, then iterated with customers using field feedback. I have also built C++ dependencies from source across CMake, Meson and Bazel projects, contributing fixes upstream to LLVM and sccache. These are different systems from your data-center hardware, but they require the same willingness to follow a problem across software layers and resolve its underlying cause.

    I take initiative when shared infrastructure can make colleagues more effective. A colleague and I started a self-hosted coding assistant that became a production service for the institute's researchers. Alongside the implementation, my work includes runbooks and onboarding documentation for our CI/CD and GPU cluster. I use coding agents for tests, reviews and debugging, while reviewing changes and running tests, linters and builds before accepting them.

    I would be keen to build on my Linux, Docker and GitLab experience with your Terraform, Ansible and Proxmox workflows. I value open technical discussions and would bring practical infrastructure ownership, curiosity and a willingness to learn the parts of your stack that are new to me. I work in English at C1 level, speak French at B1 and am open to relocating to Paris.

    \cvletterclosing{Best regards}
}
